\section*{Capítulo \ref{ch:b1:gas-exchange} --- Intercambio respiratorio de gases}

\begin{solution}{exo:b1:gas-exchange:1}
$\dot V = K(A/x)(P_1 - P_2)$: un área mayor, una barrera más delgada, una
\omterm{def:b1:gas-exchange:partial}{presión parcial} exterior más alta (ventilación) y una interior más baja
(perfusión).
\end{solution}

\begin{solution}{exo:b1:gas-exchange:2}
Treinta veces menos oxígeno por litro (\qty{7}{} frente a \qty{210}{mL/L});
ochocientas veces más densa y cincuenta veces más viscosa, y por tanto cara
de bombear; y el oxígeno difunde por ella diez mil veces más despacio, de
modo que el medio debe moverse justo por encima de la superficie.
\end{solution}

\begin{solution}{exo:b1:gas-exchange:3}
A favor de corriente: el agua y la sangre salen ambas a \qty{12}{kPa};
extracción $(21 - 12)/21 = \qty{43}{\%}$. A \omterm{def:b1:gas-exchange:gill}{contracorriente}: el agua sale a
\qty{4}{kPa} y la sangre a \qty{18}{kPa}; extracción $(21 - 4)/21 =
\qty{81}{\%}$.
\end{solution}

\begin{solution}{exo:b1:gas-exchange:4}
Disuelto (\qty{7}{\%}), unido a los grupos amino de la hemoglobina
(\qty{23}{\%}) y como bicarbonato en el \omterm{def:b1:mammal-organization:compartments}{plasma} (\qty{70}{\%}).
\end{solution}

\begin{solution}{exo:b1:gas-exchange:5}
Por minuto, \qty{6.75}{L/min}; alveolar, $15\times 0.3 = \qty{4.5}{L/min}$;
captación, $4.5\times 0.07 = \qty{0.32}{L/min}$. Superficial: por minuto
sigue siendo \qty{6.75}{L/min}, pero alveolar $30\times 0.075 =
\qty{2.25}{L/min}$ y captación \qty{0.16}{L/min}: la mitad, con el mismo
trabajo respiratorio; el espacio muerto se paga en cada respiración.
\end{solution}

\begin{solution}{exo:b1:gas-exchange:6}
$2/(6\times 0.75) = \qty{0.44}{L}$ por hora, veintidós veces su volumen de
\qty{20}{mL}.
\end{solution}

\begin{solution}{exo:b1:gas-exchange:7}
$A/x$ cae en $2\times 4 = 8$: \qty{31}{mL/min}. Para recuperar \qty{250}{},
la diferencia tendría que ser de \qty{64}{kPa}: imposible, ya que la presión
alveolar no puede superar \qty{13}{kPa} en aire y la sangre venosa no puede
bajar de cero; el paciente solo puede elevarla respirando oxígeno, y ni aun
así ocho veces.
\end{solution}

\begin{solution}{exo:b1:gas-exchange:8}
La \omterm{prop:b1:gas-exchange:transport}{anhidrasa carbónica} está dentro de la \omterm{def:b1:cell-unit-of-life:cell}{célula}, de modo que el bicarbonato
se fabrica allí y se acumula; sale por difusión a favor de su gradiente a
través de un intercambiador que admite un cloruro por cada bicarbonato, con
lo que la carga queda equilibrada. Sin ese intercambio, el bicarbonato se
quedaría dentro junto con el compañero de su protón, elevaría la osmolaridad
de la \omterm{def:b1:cell-unit-of-life:cell}{célula} y entraría agua: la \omterm{def:b1:cell-unit-of-life:cell}{célula} se hincharía en los \omterm{def:b1:body-plans-tissues:tissue}{tejidos} (y ya se
hincha un poco tal como está).
\end{solution}

\begin{solution}{exo:b1:gas-exchange:9}
El impulso de respirar procede de la subida del $\mathrm{CO_2}$. La
hiperventilación rebaja el $\mathrm{CO_2}$ muy por debajo de lo normal sin
añadir mucho oxígeno (la hemoglobina ya estaba saturada); durante la
inmersión el oxígeno cae hasta el nivel de la inconsciencia antes de que el
$\mathrm{CO_2}$ haya vuelto a subir hasta el nivel que obliga a respirar.
Sin hiperventilación, el $\mathrm{CO_2}$ alcanza ese nivel cuando todavía
sobra oxígeno, y el buceador emerge.
\end{solution}

\begin{solution}{exo:b1:gas-exchange:10}
$A = \pi(10^{-5})^2 = \qty{3.1e-10}{m^2}$; $\Delta c = JL/DA =
10^{-9}\times 10^{-3}/(2\times 10^{-5}\times 3.1\times 10^{-10}) =
\qty{160}{mol/m^3}$: mucho más que los \qty{8.7}{} disponibles, de modo que
un tubo tan estrecho no puede abastecer al músculo por difusión a
\qty{1}{mm}; hace falta un tubo más ancho (\qty{100}{\micro m}:
$\Delta c = \qty{6.4}{mol/m^3}$, justo posible) o una ventilación activa. A
\qty{10}{mm} la exigencia es diez veces mayor. Como la longitud del tubo
crece con el tamaño del cuerpo y la demanda con el volumen, la difusión por
las tráqueas limita el tamaño de los insectos a unos centímetros, y a más
solo con bombeo.
\end{solution}

\begin{solution}{exo:b1:gas-exchange:11}
Peces: agua en un sentido y sangre en el opuesto (\omterm{def:b1:gas-exchange:gill}{contracorriente}), con una
extracción del \qty{80}{\%}. Aves: aire en un sentido por los parabronquios
y sangre en ángulo recto (corriente cruzada), con una extracción intermedia
pero mayor que la de los mamíferos, y con una \omterm{prop:b1:organism-environment:surfaces}{superficie de intercambio} que
nunca recibe aire viciado. Un \omterm{def:b1:gas-exchange:lung}{pulmón} mareal mezcla el aire fresco con el gas
residual, de modo que la \omterm{prop:b1:organism-environment:surfaces}{superficie de intercambio} ve como mucho la mezcla
alveolar y la sangre solo puede equilibrarse con ella: sea cual sea la
superficie, la sangre arterial no puede superar la presión alveolar, y la
presión alveolar no puede acercarse a la inspirada mientras quede un volumen
residual.
\end{solution}

\begin{solution}{exo:b1:gas-exchange:12}
El $\mathrm{CO_2}$ es la mejor señal porque su nivel arterial responde a la
ventilación de forma directa y pronunciada (redúzcase la ventilación a la
mitad y se duplica), se mide con precisión a través del pH por los
receptores centrales, y el oxígeno, situado en la meseta de la curva de la
hemoglobina, cambia poco en el intervalo normal de ventilación: un termostato
de $\mathrm{CO_2}$ mantiene bien el oxígeno como subproducto. Falla cuando
el oxígeno cae con independencia del $\mathrm{CO_2}$, como en altitud, donde
el oxígeno inspirado es bajo pero la producción de $\mathrm{CO_2}$ no lo es,
de modo que la señal del $\mathrm{CO_2}$ dice ``respira menos'' mientras el
oxígeno dice ``respira más''; solo los cuerpos carotídeos, cuando el oxígeno
está muy bajo, la anulan, y el compromiso (hiperventilación con alcalosis)
tarda días de ajuste renal en asentarse.
\end{solution}

\begin{solution}{pb:b1:gas-exchange:1}
\textbf{1.} Aire, $210/22.4 = \qty{9.4}{mmol/L}$; agua, $7/22.4 =
\qty{0.31}{mmol/L}$.
\textbf{2.} 30.
\textbf{3.} Aire: $1.2/9.4 = \qty{0.13}{g}$; agua: $1000/0.31 =
\qty{3200}{g}$; \num{25000} veces más masa.
\textbf{4.} $210/7 = \qty{30}{L}$ de agua.
\textbf{5.} El agua es demasiado pesada y demasiado pobre en oxígeno para
empujarla hacia dentro y hacia fuera; hay que hacerla pasar una sola vez, lo
que además hace posible la \omterm{def:b1:gas-exchange:gill}{contracorriente}; el aire es lo bastante barato
para moverlo dos veces.
\textbf{6.} Por minuto, $12\times 0.5 = \qty{6}{L/min}$; alveolar,
$12\times 0.35 = \qty{4.2}{L/min}$.
\textbf{7.} $150/500 = \qty{30}{\%}$ de cada respiración, y de la ventilación por minuto.
\textbf{8.} Entregado, $4.2\times 0.21 = \qty{0.88}{L/min}$; captado,
$4.2\times 0.07 = \qty{0.29}{L/min}$, cercano a los \qty{250}{mL/min}.
\textbf{9.} $0.25/(6\times 0.21) = \qty{20}{\%}$ (el \qty{28}{\%} de lo que
llega a los \omterm{def:b1:gas-exchange:lung}{alvéolos}).
\textbf{10.} $0.14\times(101 - 6.3) = \qty{13.3}{kPa}$.
\textbf{11.} $250/5 = \qty{50}{mL/L}$; venoso \qty{150}{mL/L}, saturado al
\qty{75}{\%}.
\textbf{12.} $3000/20 = \qty{150}{mL/L}$; venoso \qty{50}{mL/L}, saturado al
\qty{25}{\%}.
\textbf{13.} $50/60 = \qty{0.83}{mL/min}$.
\textbf{14.} $0.83/(7\times 0.8) = \qty{0.149}{L/min}$, \qty{8.9}{L} por
hora.
\textbf{15.} $0.83/(7\times 0.43) = \qty{0.28}{L/min}$: la \omterm{def:b1:gas-exchange:gill}{contracorriente}
reduce a la mitad el agua que hay que bombear, y el trabajo.
\textbf{16.} \qty{8.9}{kg} de agua por hora; el ser humano mueve $6\times
60\times 1.2 = \qty{430}{g}$ de aire por hora --- una veinteava parte de
la masa, para trescientas veces el oxígeno ($15\,000$ frente a
\qty{50}{mL} por hora).
\textbf{17.} $100\times(0.95 - 0.15) = \qty{80}{mL/L}$; gasto cardíaco,
$0.83/0.080 = \qty{10.4}{mL/min}$.
\textbf{18.} Agua/sangre, $149/10.4 = 14$; aire/sangre, $4.2/5 = 0.84$. El
pez debe hacer pasar catorce volúmenes de agua por cada volumen de sangre
porque cada uno contiene muy poco oxígeno; el ser humano, alrededor de uno.
\textbf{19.} Captación $\times 1.5$ de un agua que contiene $5.5/7 = 0.79$
veces lo de antes: ventilación $\times 1.9$; el coste de respirar,
aproximadamente proporcional al caudal (o más, ya que sube más deprisa que
linealmente), pasa del \qty{10}{\%} a por lo menos el \qty{13}{\%} de una
captación mayor: el pez trabaja más para respirar en agua caliente.
\textbf{20.} Extremo de entrada del agua: $21 - 18 = \qty{3}{kPa}$; extremo
de salida del agua: $4 - 3 = \qty{1}{kPa}$: positiva en ambos extremos.
\textbf{21.} Extremo de entrada, $21 - 3 = \qty{18}{kPa}$; extremo de
salida, $12 - 12 =
0$: los dos fluidos se han equilibrado y no queda
gradiente, aunque el agua todavía conserva el \qty{57}{\%} de su oxígeno.
\textbf{22.} A la de favor de corriente: la sangre se equilibra con un único
gas alveolar bien mezclado y sale a su presión; no puede superar los
\qty{13.3}{kPa} porque aguas abajo no hay gas más fresco al que acercarse,
como sí lo hay en la \omterm{def:b1:gas-exchange:gill}{contracorriente}.
\textbf{23.} Trucha, $2000/(2\times 10^{-4}) = \num{1e7}$ (cm$^2$ por cm);
ser humano, $10^6/(5\times 10^{-5}) = \num{2e10}$: cociente \num{2000}.
Captaciones, $250/0.83 = 300$. La trucha toma siete veces más oxígeno por
unidad de $A/x$: su diferencia media de presión a través de la laminilla
debe ser mayor, y eso es lo que aporta la \omterm{def:b1:gas-exchange:gill}{contracorriente}, que mantiene un
gradiente a lo largo de toda la superficie donde el del \omterm{def:b1:gas-exchange:lung}{pulmón} mareal es
pequeño cerca del final del equilibrado.
\textbf{24.} El $\mathrm{CO_2}$ es veinticinco veces más soluble en agua que
el oxígeno: el agua se lo lleva tan deprisa como se produce, y el
$\mathrm{CO_2}$ de la sangre de un pez se mantiene muy bajo.
\textbf{25.} La \omterm{def:b1:gas-exchange:gill}{branquia}, alrededor del \qty{80}{\%}, y el \omterm{def:b1:gas-exchange:lung}{pulmón}, alrededor
del \qty{20}{\%} del oxígeno del medio inspirado; la diferencia la marcan el
flujo en un solo sentido de la \omterm{def:b1:gas-exchange:gill}{branquia} y su disposición a \omterm{def:b1:gas-exchange:gill}{contracorriente}
del agua y la sangre.
\end{solution}
